\documentclass[12pt]{article}
\usepackage{graphicx} % Required for inserting images
\usepackage[margin=0.5in]{geometry}
\usepackage{amsmath, amssymb}
\usepackage{mathrsfs}


\usepackage{soul}


% Dr. David Brown Commands
\newcommand{\ds}{\displaystyle}
\newcommand{\R}{\mathbb{R}}
\newcommand{\N}{\mathbb{N}}
\newcommand{\Q}{\mathbb{Q}}
\newcommand{\C}{\mathbb{C}}


% Commands to get script r
\usepackage{calligra}
\DeclareMathAlphabet{\mathcalligra}{T1}{calligra}{m}{n}
\DeclareFontShape{T1}{calligra}{m}{n}{<->s*[2.2]callig15}{}
\newcommand{\scripty}[1]{\ensuremath{\mathcalligra{#1}}}

% Unit commands
\newcommand{\degree}{\text{°}}
\newcommand{\unitSpeed}{\frac{\text{m}}{\text{s}}}
\newcommand{\unitAcceleration}{\frac{\text{m}}{\text{s}^2}}

% Constant commands
\newcommand{\avogadro}{6.022\times10^{23}}

\newcommand{\boltzmannConstK}{1.381\times10^{-23}\frac{\text{J}}{\text{K}}}
\newcommand{\boltzmannConstInEV}{8.617\times10^{-5}\frac{\text{eV}}{\text{K}}}
\newcommand{\planckConst}{6.626\times10^{-34}\text{J}\cdot\text{s}}

\newcommand{\atmP}{1.01325\times10^5\,\text{Pa}}

% Known Value Commands
\newcommand{\electronMass}{9.109\times10^{-31}\,\text{kg}}
\newcommand{\protonMass}{1.673\times10^{-27}\,\text{kg}}

% Commands to easily write partial derivatives
\newcommand{\partDeriv}[2]{\frac{\partial #1}{\partial #2}}
\newcommand{\doublePartDeriv}[2]{\frac{\partial^2 #1}{\partial^2 #2}}


\title{Problem 7.26}
\author{Gabriel Decker}


\begin{document}


\maketitle
\begin{flushleft}


In this problem, we consider ${}^3\text{He}$, modeling it as a nearly degenerate Fermi gas (I say nearly degenerate because $T$ is nearly zero, but not zero).
We're given the fact that Helium-3 atoms are fermions because of the unpaired neutrons in their nuclei.
We're also ignoring the fact that ${}^3\text{He}$ liquifies.\\~\\

\textbf{Part a}\\
First, we'll find the Fermi energy and temperature.
These are given by the following equations.
\begin{equation}
    \epsilon_F=\frac{h^2}{8m}\left(\frac{3N}{\pi V}\right)^{2/3}
\end{equation}
\begin{equation}
    T_F=\epsilon_F/k
\end{equation}\\~\\

We are given that the molar volume (the volume of one mole of ${}^3\text{He}$) is 37 cm$^3$.
Or, $N/V=\avogadro/(37\times10^{-6}\,\text{m}^3)=1.62757\times10^{28}\,\text{m}^{-3}$.\\~\\

Now, let's plug in our data into the above definitions to get our desired quantities.
I did this in a python program.
$$\epsilon_F=6.82\times10^{-23}\,\text{J}$$
$$T_F=4.94\,\text{K}$$\\~\\~\\




\textbf{Part b}\\
We will now use the above to calculate $C_V$, assuming temperatures $T\ll T_F$.
We will compare our result with the experimental result of $C_V=(2.8\,\text{K}^{-1})NkT$.
The equation for heat capacity in this state is the following.
\begin{equation}
    C_V=\frac{\pi^2 Nk^2 T}{2\epsilon_F}
\end{equation}\\~\\

We can simplify this to get it into the right form and then calculate the result.
$$C_V=\frac{\pi^2 Nk^2 T}{2\epsilon_F}$$
$$\implies C_V=\frac{\pi^2k}{2\epsilon_F}NkT$$
$$\implies C_V=(0.999\,\text{K}^{-1})NkT$$
The experimental result is about three times larger.\\~\\~\\





\textbf{Part c}\\
Now, we will compare this nearly degenerate Fermi gas model of ${}^3\text{He}$ to solid ${}^3\text{He}$ by comparing their entropies.
We're given that for solid ${}^3\text{He}$, its entropy depends nearly entirely on the nuclear spin arrangements.
This implies that for each nucleus, there are two states, and that the entropy in this case is $S_s=k\ln(\Omega)=k\ln(2^N)$, or $S_s=Nk\ln(2)$.\\~\\

For our nearly degenerate Fermi gas model, we already have an equation for the heat capacity, so we can use the following equation to get the entropy.
\begin{equation}
    dS=\frac{C_V}{T}\,dT
\end{equation}
We'll use our result from part b and integrate this.
$$S_g=\int_0^T\frac{C_V}{T'}\,dT'$$
$$\implies S_g=\int_0^T\frac{(0.999\,\text{K}^{-1})NkT'}{T'}\,dT'$$
$$\implies S_g=(0.999\,\text{K}^{-1})Nk\int_0^T\,dT'$$
$$\implies S_g=(0.999\,\text{K}^{-1})NkT$$\\~\\

The following is a plot of both entropies as a function of $T$ for $N=1$.\\
\includegraphics[width=0.6\linewidth]{entropies.png}\\~\\

We can also find at what temperature these entropies are equal.
$$(0.999\,\text{K}^{-1})NkT=Nk\ln(2)$$
$$\implies T=\frac{\ln(2)}{0.999}\,\text{K}$$
$$\implies T=0.694\,\text{K}$$






\end{flushleft}

\end{document}
